\section*{Appendix B Supplementary Tables and Figures}
\addcontentsline{toc}{section}{Appendix B Supplementary Tables and Figures}
\setcounter{subsection}{0}
\renewcommand{\thesubsection}{B.\arabic{subsection}}
\setcounter{figure}{0}
\renewcommand{\thefigure}{B.\arabic{figure}}
\setcounter{table}{0}
\renewcommand{\thetable}{B.\arabic{table}}

\subsection{Descriptor construction and comparison objectives}\label{app:descriptors}
Table~\ref{tab:descriptors} describes the information available from composition alone. Elemental tables use fixed source conventions from Mendeleev 1.3.0, Magpie 0.9.3 and tabulated Zunger radii. Material-level physical responses enter the target and class definitions, not these inputs. Source values that are zero remain zero; unavailable entries retain their missingness indicators.
\par
\begin{table}[H]
\centering
\caption{Composition-derived information used by the surrogate}\label{tab:descriptors}
\begin{tabularx}{\linewidth}{@{}l X@{}}
\toprule
Block & Construction and interpretation\\
\midrule
A and B sites & Separate elemental descriptors; signed and absolute A--B property differences. Site order is retained.\\
X composition & O/N/F/S counts, seven family indicators and formal group charge. X is a three-atom group.\\
X statistics & Stoichiometric means, standard deviations, minima and maxima; A--X and B--X contrasts.\\
Electronic information & Valence/unfilled $s,p,d,f$ counts, Pettifor number, Martynov--Batsanov electronegativity and Zunger $r_s+r_p$.\\
Geometry & Charge-pair counts, radii, tolerance and octahedral factors, and $\tau$ statistics over admissible assignments. Strict and approximate quantities remain distinct.\\
Missingness & Indicators accompany quantities throughout. Tree fitting handles unavailable entries; baseline median imputation uses fitting data only.\\
\bottomrule
\end{tabularx}
\end{table}
\par

The radius construction distinguishes three operations:
\begin{itemize}
  \item \textbf{Charge balance:} enumerate positive A/B charges balancing O$^{2-}$, N$^{3-}$, F$^{-}$ and S$^{2-}$ in X.
  \item \textbf{Effective X radius:} take the stoichiometric mean of O/N/F/S radii 1.40/1.46/1.33/1.84~\AA.
  \item \textbf{Coordination:} require A-XII and B-VI for strict values. The approximation permits A-VIII or A-VI when XII is unavailable. N-containing groups have approximate, flagged values only.
\end{itemize}
The geometric measures over admissible assignments are
\begin{samepage}
\begin{equation}
 t=\frac{r_A+r_X}{\sqrt{2}(r_B+r_X)},\qquad
 \mu=\frac{r_B}{r_X},\qquad
 \tau=\frac{r_X}{r_B}-q_A\left(q_A-\frac{r_A/r_B}{\ln(r_A/r_B)}\right).
 \label{eq:geometry}
\end{equation}
\noindent Here $t$ is the tolerance factor, $\mu$ the octahedral radius ratio and $\tau$ the alternative tolerance descriptor. The assigned radii are $r_A,r_B,r_X$, and $q_A$ is the positive formal A-site charge. The $\tau$ expression requires $r_A/r_B>1+10^{-6}$. Summaries over assignments do not establish a unique valence state or mixed-anion stability.
\end{samepage}

Table~\ref{tab:routes} distinguishes numerical objectives from screening objectives. All routes use the same outer grouped validation. For probability and ranking models, calibration is confined to the corresponding fitting data. The single-query ranker uses the LightGBM 3.3.5 implementation with a query-size capacity sufficient for the entire fitting pool.
\par
\begin{table}[H]
\centering
\caption{Compared objectives and prediction scales}\label{tab:routes}
\begin{tabularx}{\linewidth}{@{}l X@{}}
\toprule
Method & Objective and output\\
\midrule
R-raw & Squared loss on $V$; negative predictions clipped to zero.\\
R-log & Squared loss on $\log(1+V)$; exponentiate, subtract one and clip at zero. This transformation need not estimate the original-scale conditional mean.\\
W-log & Weighted R-log. Log-target bins have width 0.2; density is smoothed over five bins with a Gaussian kernel. Inverse-square-root density weights are capped at 10 before mean normalization.\\
H-soft & Binary nonmetal probability and two conditional squared-loss regressions; Eq.~\eqref{eq:mixture} returns original-scale $\widehat V$.\\
S-eps & Product of nonmetal and conditional near-positive probabilities, Eq.~\eqref{eq:chain}.\\
B-eps & Binary loss for $V\le2$ with inverse class-frequency weights, followed by probability calibration.\\
L-single / L-hash4 & LambdaRank with binary relevance $\mathbf{1}\{V\le2\}$, gains 0/1 and truncation level 23; one query or four A--B hash queries.\\
Additive Ridge & A/B/X one-hot descriptors and Ridge penalty 1; no interaction terms.\\
Mean / random & Training-fold mean of $V$; or a fixed hash ordering independent of the response.\\
\bottomrule
\end{tabularx}
\end{table}
\par
Probability calibration uses a logistic map $\sigma(a+bz)$, where $z$ is the raw score normalized by the fitting model's score mean and standard deviation, and $\sigma(t)=(1+e^{-t})^{-1}$. L-BFGS-B minimizes binomial negative log-likelihood plus $b^2/2$, with $a\in[-35,35]$ and $b\in[0,50]$. A single-class sample uses $(n_++1)/(n+2)$, where $n_+$ of $n$ fitting records satisfy the event.
\FloatBarrier

\subsection{Validation counts and residual diagnostics}\label{app:validation}
Table~\ref{tab:folds} retains the denominators behind the main enrichment result. Under a uniform sample without replacement, the probability of at least the observed $h_k$ near-positive cases is
\begin{samepage}
\begin{equation}
 P(H\ge h_k\mid N,M,k)=
 \sum_{j=h_k}^{\min(k,M)}
 \frac{\binom{M}{j}\binom{N-M}{k-j}}{\binom{N}{k}}.
 \label{eq:hypergeom}
\end{equation}
\noindent Here $H$ is the random hit count under uniform selection, and $N,M,k$ have the pool-specific meanings in Eq.~\eqref{eq:ef}. Impossible binomial terms are zero. Holm correction is applied within the specified family of comparisons; the two file-03 candidate lists form their own family. The test compares a list with uniform sampling, not with all possible learned selectors. Because model development precedes this comparison, its small tail probability is descriptive evidence of enrichment within the evaluated pool.
\end{samepage}

\input{tables/fold_counts}
Table~\ref{tab:nested} evaluates selection within the outer training partitions. Numerical selection chooses H-soft in all five folds. Screening selection chooses among the candidate routes within each training partition and is evaluated on that partition's untouched outer fold.
\input{tables/nested}
\begin{samepage}
For file 03, $N=3{,}969$, $M=41$, $k=20$ and $h_k=13$ give a tail probability of $1.37\times10^{-22}$ (two-comparison Holm: $2.73\times10^{-22}$). This tests $V\le2$; the hidden test responses provide no observed hit count.\par
\end{samepage}

Figure~\ref{fig:residuals} separates mean signed error from error magnitude for each X family. The heat strip encodes the mean residual $\widehat V-V$ with blue for negative and amber for positive values. In panel (b), solid bars denote RMSE and hatched bars the empirical 95th percentile of absolute residuals. Each family has 1,608 outer-fold predictions. Small mean bias can coexist with substantial individual errors, particularly for $\mathrm{N_3}$ and $\mathrm{N_2O}$.
\paperfigure{A1/A1_residual_structure_v3.pdf}{Residual bias and error scale across X families}{fig:residuals}
\FloatBarrier

\subsection{Probability calibration and query assignment}
Figure~\ref{fig:calibration} compares predicted probabilities with observed near-positive frequencies using outer-fold predictions. All four panels concern $V\le2$. The diagonal denotes agreement; the slope labels summarize separate calibration diagnostics. A reliability improvement changes probability interpretation and need not improve the within-fold ranking.
\paperfigure{A2/A2_calibration_v2.pdf}{Probability reliability on held-out A--B groups}{fig:calibration}

Table~\ref{tab:calibration} gives the pooled diagnostic coefficients. Calibration changes the mean nonmetal Brier score only slightly, from 0.015913 to 0.015893; the conditional-event score changes from 0.156119 to 0.155870. The B-eps score changes from 0.007372 to 0.005750. The S-eps product remains imperfectly calibrated, with diagnostic slope 1.7617. These checks motivate keeping ranking quality and probability accuracy as distinct results.
\input{tables/calibration}

For B-eps and both ranking arrangements, each outer-fold Top-20 list is unchanged by calibration. However, combining predictions across folds yields raw/calibrated overlaps of only 15 or 16 out of 20, because their scales differ. The primary endpoint therefore remains the mean of within-fold enrichments. Table~\ref{tab:query} also shows the effect of changing the four-query assignment: about 15 of 20 selections remain shared on average. These comparisons use the same outer A--B partitions and fixed model capacity.
\input{tables/query}
\FloatBarrier

\subsection{Prediction-table stability}
The four selection pools provide complementary checks of how prediction variation changes decisions. Pools 07 and 11 fix X at $\mathrm{NO_2}$ and contain 2,704 and 529 candidates. Pools 08 and 12 allow $\mathrm{O_3}$, $\mathrm{NO_2}$ and $\mathrm{FO_2}$ and contain 8,112 and 1,587 candidates. These calculations use the file-02 surrogate tables to assess the P1 output, without solving the later optimization problems.

Let $v_0(c)$ denote the nominal table, $v_b(c)$ bootstrap-member $b$, and $c_0$ and $c_b$ their respective minimizers over the same allowed pool. The two directions of decision cost are
\begin{samepage}
\begin{equation}
 R_b=v_b(c_0)-v_b(c_b),\qquad
 \widetilde R_b=v_0(c_b)-v_0(c_0).
 \label{eq:surrogate-regret}
\end{equation}
\noindent Here $R_b$ measures the nominal choice under member $b$, while $\widetilde R_b$ measures the member's choice under the nominal table. Both are nonnegative and use predicted violation degrees. They must not be interpreted as regret against measured material properties.
\end{samepage}

Figure~\ref{fig:stability} combines rank distributions with these costs. A low rank frequency alone does not show that the nominal candidate performs poorly, so the right panel reports regret quantiles as well. Table~\ref{tab:nominal} preserves the numerical ranges. Neither direction is uniformly negligible: for example, pool 12 has a reverse-regret 95th percentile of 6.6253. None of the four pools satisfies the specified criterion that both maximum regrets stay below 0.1.
\paperfigure{A3/A3_decision_stability_v2.pdf}{Nominal-candidate ranks and bidirectional surrogate regret}{fig:stability}
\input{tables/nominal_stability}
\FloatBarrier

Figure~\ref{fig:splithalf} checks whether averaging makes the decision tables more reproducible. Each of 1,000 fixed partitions divides the same 20 bootstrap models into two groups of ten. Their mean tables are compared with all cross-group individual-member pairs. The left panel reports agreement on the minimum; the right compares symmetric cross-table regret. Table~\ref{tab:stability} adds the Top-20 overlap. These 1,000 partitions reuse 20 models and are not 1,000 independent fitting experiments.
\paperfigure{A4/A4_split_half_v2.pdf}{Split-half effects of averaging on decision stability}{fig:splithalf}
\input{tables/stability}
Averaging improves shortlist agreement across all four pools and reduces the mean cross-table cost. Nevertheless, the minimum candidate is much less reproducible in pools 11 and 12 than in pools 07 and 08. This supports reporting a shortlist and prediction variation alongside the nominal minimum, without treating table agreement as a measure of physical accuracy.
\FloatBarrier

\subsection{Extrapolation protocols}
Table~\ref{tab:stress} reports three exclusion designs:
\begin{itemize}
  \item \textbf{Element:} remove selected elements from fitting and test their compositions. The five evaluations can contain overlapping materials.
  \item \textbf{Target region:} withhold A--B groups containing records below 0.5, 1 or 2; add background groups to reach the specified 20\% test share.
  \item \textbf{X family:} cross seven families with five A--B folds. Fitting excludes both the tested family and the tested A--B groups.
\end{itemize}
\input{tables/stress}

In Figure~\ref{fig:stress}, large labels are the recorded metric values. Smaller labels and color express percentage changes from each route's own grouped-validation reference: R-log RMSE 4.9048, R-log EF 60.4863 and S-eps EF 62.2234. Blue and amber encode negative and positive changes, respectively. Lower RMSE and higher EF are desirable, so one color does not have a common good/bad meaning across columns.
\paperfigure[0.92\linewidth]{A5/A5_stress_protocols_v4.pdf}{Exclusion-protocol outcomes: protocol-level evidence, not final-model validation}{fig:stress}

The X-family protocol gives R-log RMSE 6.5035 and mean EF 12.9015 over 29 valid enrichment evaluations. The six zero-positive pools retain numerical-error results but have undefined enrichment. Target-region exclusion at $V<2$ removes the positive fitting examples required by the S-eps conditional event model; its fallback is the smoothed prior, leaving the nonmetal probability to determine ordering. This case is not a demonstration of learning an unseen zero-violation class.
\FloatBarrier

\subsection{Descriptor and ablation comparisons}
Table~\ref{tab:features} separates a descriptor ladder under R-log from reverse ablations of H-soft. The ladder measures the benefit of broader composition information under one regression objective. The reverse ablations remove complete dependency blocks from the final model, including associated statistics and missing-value indicators. Because some geometric and X-derived descriptors overlap, the resulting effects are not additive.
\input{tables/features}

Figure~\ref{fig:ablation} displays reverse-ablation differences from full H-soft. A marker locates the paired point estimate, horizontal whiskers give the 95\% bootstrap interval, and the vertical reference denotes zero effect. An interval crossing zero is shown in neutral gray. The RMSE and EF panels have different desirable directions, as indicated by their signed changes.
\paperfigure{A6/A6_ablation_v2.pdf}{Reverse-ablation effects with paired bootstrap intervals}{fig:ablation}

Table~\ref{tab:paired} gives the complete paired comparisons used in the main text. Resampling is performed within each outer A--B fold, with the same multiplicities shared by compared methods. All 5,000 resamples give a defined five-fold enrichment mean. The intervals are conditional and unadjusted for multiple interval comparisons; they support the stated comparisons rather than a universal ranking of models or descriptor families.
\input{tables/paired}
\FloatBarrier
